\documentclass[10pt]{article}
\usepackage[margin=26mm]{geometry}
\usepackage{amsmath,amssymb,amsthm}
\usepackage[hidelinks]{hyperref}
\newtheorem{proposition}{Proposition}
\newtheorem{corollary}[proposition]{Corollary}
\newcommand{\HK}{\operatorname{HK}}
\newcommand{\HH}{\operatorname{HH}}
\newcommand{\im}{\operatorname{im}}
\hypersetup{pdftitle={A characteristic obstruction in higher Koszul homology},pdfauthor={Alec Kriebel}}
\title{\vspace{-10mm}A characteristic obstruction in\\ higher Koszul homology}
\author{Alec Kriebel\thanks{OpenAI tools were used extensively in development,
 solving, drafting, coding and adversarial verification. This is an unrefereed
 preprint; no completed external human peer review or formal proof-assistant
 verification is claimed.}\\
 \small Independent researcher\\
 \small\href{https://orcid.org/0009-0001-9320-500X}{ORCID: 0009-0001-9320-500X}}
\date{October 14, 2026 \quad Version 1.0.0}
\begin{document}
\maketitle
\begin{abstract}
Over every field of positive characteristic, the free algebra on two
 degree-one generators is Koszul but has nonzero positive-degree higher
 Koszul homology. We give an explicit class and the complete cyclic-word
 calculation. This elementary application of classical tensor-algebra
 formulas refutes the unrestricted-field equivalence reported in
 Oberwolfach Report 10/2016. The characteristic-zero converse remains
 unresolved here.
\end{abstract}

\section{The unrestricted statement and its scope}
Solotar's contribution to Oberwolfach Report 10/2016
\cite[printed pp.~479--481]{OWR} asks whether a quadratic algebra is
Koszul if and only if its higher Koszul homology vanishes in positive
degrees. The contribution states no restriction on the characteristic.
We use higher homology with coefficients in the algebra itself and
fundamental cocycle \(e_A\), as defined by Berger--Lambre--Solotar
\cite[\S5.2]{BLS}. Read over unrestricted fields, the equivalence is false.

This qualification matters. The forward theorem in
\cite[Theorem~6.4]{BLS} assumes characteristic zero; the adjacent
Conjecture~6.5 includes the additional condition \(\HK^{\mathrm{hi}}_0(A)\cong k\)
and is motivated as a converse, while its preceding discussion also
considers removing the characteristic hypothesis.
Our counterexample refutes the unrestricted forward implication of
either formulation. It neither proves nor disproves the
characteristic-zero vanishing-to-Koszulness converse, and does not
establish a unique historical intended domain.

The calculation is an elementary assembled consequence of the classical
tensor-algebra Hochschild formulas of Loday--Quillen \cite[\S5]{LQ},
Goodwillie's Euler cap operator \cite[II.4.8]{G}, and the BLS definition.
No new homological mechanism or historical firstness is claimed.
We record the explicit obstruction to the unrestricted written statement.

\section{The free-algebra calculation}
Let \(k\) be a field, let \(V\) have a finite basis of \(r\) letters,
and put \(A=T_k(V)\), graded by word length. This is a quadratic
algebra with relation space \(R=0\).
Write \(\HK_i^{\mathrm{hi}}(A)[n]\) for total internal weight \(n\):
a chain in \(A_{n-i}\otimes W_i\) has that total weight.
For \(n>0\), let \(\mathcal O_n\) denote the cyclic orbits of length-\(n\)
words in the chosen letters. Orbit cardinality means the number of
\emph{distinct} rotations, which can be smaller than \(n\).

\begin{proposition}[Classical specialization]\label{prop:orbits}
The algebra \(A\) is Koszul over every field. If \(\operatorname{char}k=p>0\),
then, for \(n>0\) and \(i=0,1\),
\[
 \HK_i^{\mathrm{hi}}(A)[n]
 \cong\bigoplus_{\substack{O\in\mathcal O_n\\p\mid |O|}}k.
\]
In characteristic zero these positive-weight groups vanish.
For all fields, \(\HK_i^{\mathrm{hi}}(A)=0\) for \(i\ge2\),
\(\HK_0^{\mathrm{hi}}(A)[0]=k\), and
\(\HK_1^{\mathrm{hi}}(A)[0]=0\).
\end{proposition}

\begin{proof}
Multiplication \(\mu(a\otimes v)=av\) maps the word/letter basis
of \(A\otimes V\) bijectively onto the nonempty words of \(A\).
Consequently
\[
 0\longrightarrow A\otimes V\xrightarrow{\ \mu\ }A
 \xrightarrow{\ \varepsilon\ }k\longrightarrow0
\]
is an exact linear graded free left resolution, proving Koszulness.
Here \(\varepsilon\) is augmentation. This map \(\mu\) is distinct
from the coefficient-\(A\) commutator differential used next.

The quadratic Koszul spaces satisfy \(W_0=k\), \(W_1=V\), and
\(W_j=0\) for \(j\ge2\), since every defining intersection contains
\(R=0\). Thus the ordinary coefficient complex is exactly
\[
 0\longrightarrow A\otimes V\xrightarrow{\ b\ }A\longrightarrow0,
 \qquad b(a\otimes v)=av-va.
\]
In particular \(\HK_1(A)=\ker b\), \(\HK_0(A)=\operatorname{coker}b\),
and \(\HK_j(A)=0\) for \(j\ge2\).
The fundamental left cap product on \(a\otimes v\) is \(va\).
For a cycle \(z=\sum a_j\otimes v_j\), the equation \(b(z)=0\)
gives \(\sum v_ja_j=\sum a_jv_j\). Hence the induced higher
differential is
\[
 \delta:\ker b\longrightarrow\operatorname{coker}b,
 \qquad \delta[z]=\Big[\sum_j a_jv_j\Big].
\]
This definition is valid also in characteristic two
\cite[\S5.2]{BLS}. There is no incoming higher boundary from \(\HK_2(A)\).

Fix \(n>0\). Concatenation identifies \(A_{n-1}\otimes V\)
and \(A_n\) separately with the vector space \(E_n\) on length-\(n\)
words. Let \(\rho\) move the final letter to the beginning.
Then \(b=1-\rho\); this is the classical small complex
\cite[Lemmas~5.1--5.2]{LQ}. The higher map becomes the natural map
\[
 E_n^\rho\longrightarrow (E_n)_\rho
       :=E_n/(1-\rho)E_n.
\]
On an orbit \(O\) of cardinality \(d\), equality of neighboring
coefficients makes the invariant subspace the line generated by
\(s_O=\sum_{w\in O}w\). The coinvariant quotient is the line
generated by any \([w]\), which is nonzero: the coefficient-sum
functional kills all differences and takes \([w]\) to one.
No division by \(d\) is involved. The induced map is
\[
 s_O\longmapsto d[w].
\]
Its kernel and cokernel are both one-dimensional precisely when
\(d=0\) in \(k\), and otherwise both vanish. Summing the orbit
blocks proves the formula. At weight zero only the degree-zero
chain \(k\) occurs; all assertions in degrees at least two follow
from the absence of their ordinary homology groups.
\end{proof}

\begin{corollary}\label{cor:counterexample}
For every field \(k\) of positive characteristic \(p\), the quadratic
algebra \(k\langle x,y\rangle\) is Koszul and
\(\HK_1^{\mathrm{hi}}(k\langle x,y\rangle)[p]\ne0\).
\end{corollary}
\begin{proof}
The word \(w=x^{p-1}y\) has exactly \(p\) distinct rotations
\(w_j=\rho^j(w)\), because the unique \(y\) changes position.
Write \(w_j=a_jv_j\), with \(v_j\) its final letter. The chain
\[
 z_p=\sum_{j=0}^{p-1}a_j\otimes v_j
\]
is a sum of distinct basis tensors and is therefore nonzero.
Its commutator boundary telescopes, and
\(\delta[z_p]=p[w]=0\). There are neither ordinary boundaries
entering degree one nor higher boundaries entering from \(\HK_2\).
Thus \([z_p]\ne0\) in higher degree one.
In characteristic two, \(z_2=x\otimes y+y\otimes x\).
\end{proof}

The same orbit contributes a positive-weight higher degree-zero class,
so adding \(\HK_0^{\mathrm{hi}}(A)\cong k\) cannot rescue the unrestricted
equivalence. In contrast, one generator has only singleton orbits;
its positive-weight higher groups vanish in every characteristic.
In particular, listing repeated rotations of \(x^p\) cannot supply
a counterexample.

\section{Classical operators, verification and priority limits}
The operator here is Euler cap, not Connes' differential \(B\).
Loday--Quillen \cite[Lemma~5.3]{LQ} identify \(B\) with the cyclic
norm in the opposite direction. On an orbit of size \(d\) in
weight \(n\), Euler cap has scalar \(d\), whereas \(B\) has scalar
\(n/d\). Their composites multiply by \(n\), consistently with
Goodwillie's identity \cite[Theorem~II.4.2]{G}. For the primitive
weight-\(p\) orbit above, \(B\) is invertible in characteristic \(p\)
while Euler cap is zero. The two homologies cannot be substituted
for each other merely from characteristic dependence.

The supporting archive includes editable source, all 65 frozen original
files and four evidence archives, independent algebra and source-scope
reviews, and a dated bounded priority audit. Fresh exact controls compare
234 matrix/necklace cases, check 14 prime witnesses and ten negative
controls, and run 2110 explicit checks in both normal and optimized
Python. Corrupted expectations are rejected in both modes.
These finite and integrity checks supplement the unlimited algebraic
proof; they do not certify it or determine historical priority.

No earlier explicitly stated exact counterexample was located in the
inspected corpus. This bounded finding does not establish worldwide
firstness or current openness. The entire mechanism receives the
classical credit stated above. Version and access limits, including
unavailable complete journal-final bodies for some related work,
are recorded in the accompanying audit. Downloaded primary publications
are excluded from the deposit; bibliographic locators and hashes remain.

\begin{thebibliography}{9}
\bibitem{OWR}
A.~Solotar, Koszul Calculus (joint work with R.~Berger and T.~Lambre),
in \emph{Hochschild Cohomology in Algebra, Geometry, and Topology},
Oberwolfach Reports \textbf{13} (2016), 449--506,
contribution pp.~479--481.
\href{https://doi.org/10.4171/OWR/2016/10}{doi:10.4171/OWR/2016/10}.
\bibitem{BLS}
R.~Berger, T.~Lambre and A.~Solotar, Koszul calculus,
\emph{Glasgow Math. J.} \textbf{60} (2018), 361--399.
\href{https://doi.org/10.1017/S0017089517000167}{doi:10.1017/S0017089517000167}.
First preprint: \href{https://arxiv.org/abs/1512.00183v1}{arXiv:1512.00183v1}
(December 1, 2015).
\bibitem{LQ}
J.-L.~Loday and D.~Quillen,
Cyclic homology and the Lie algebra homology of matrices,
\emph{Comment. Math. Helv.} \textbf{59} (1984), 565--591.
\href{https://doi.org/10.1007/BF02566367}{doi:10.1007/BF02566367}.
\bibitem{G}
T.~G.~Goodwillie, Cyclic homology, derivations, and the free loopspace,
\emph{Topology} \textbf{24} (1985), 187--215.
\href{https://doi.org/10.1016/0040-9383(85)90055-2}{doi:10.1016/0040-9383(85)90055-2}.
\end{thebibliography}
\end{document}
